\section{Introduction}
\label{sec:introduction}

For finite simple graphs $H$ and $J$, the Ramsey number $R(H,J)$ is the
least integer $N$ such that every red--blue colouring of the edges of the
complete graph on $N$ vertices contains a red copy of $H$ or a blue copy of
$J$. Copies need not be induced. Write $C_m$ for the cycle on exactly $m$
vertices and $K_n$ for the complete graph on $n$ vertices.

\begin{theorem}\label{thm:main}
For all integers $m\ge n\ge3$ with $(m,n)\ne(3,3)$,
\[
 R(C_m,K_n)=(m-1)(n-1)+1.
\]
For the excluded pair, $R(C_3,K_3)=6$.
\end{theorem}

Theorem~\ref{thm:main} establishes the cycle--clique conjecture of
Erd\H{o}s, Faudree, Rousseau and Schelp
\cite{ErdosFaudreeRousseauSchelp1978}, in the formulation stated by
Keevash, Long and Skokan~\cite{KeevashLongSkokan2021}.
The cycle length in the theorem is exact. Its lower bound comes from
$n-1$ disjoint red cliques of order $m-1$, with all edges between them blue.
The task is to show that one additional vertex forces one of the two
specified graphs.

\subsection{Previous results and the parameter range}

The case $n=3$ goes back to Chartrand and Schuster
\cite{ChartrandSchuster1971}; see also Bondy and Erd\H{o}s
\cite[p.~47]{BondyErdos1973}. Bondy and Erd\H{o}s
\cite[Theorem~4]{BondyErdos1973} established the formula when
$m\ge n^2-2$. The complete conjectured range was established for $n=4$
by Yang, Huang and Zhang~\cite{YangHuangZhang1999}, for $n=5$ by
Bollob\'as and coauthors~\cite{BollobasEtAl2000}, for $n=6$ by
Schiermeyer~\cite{Schiermeyer2003}, and for $n=7$ by Chen, Cheng and
Zhang~\cite{ChenChengZhang2008}.
For $n=8$, further cases were established: $m=8$
\cite{JaradatAlzaleq2007,ZhangZhang2009}, $m=9$
\cite{BatainehJaradatAlZaleq2011}, and $10\le m\le15$
\cite{Baniabedalruhman2023}.
Theorem~1 of Nikiforov~\cite{Nikiforov2005} states the formula for
$n\ge4$ and $m\ge4n+2$.

Keevash, Long and Skokan~\cite[Theorem~1.1]{KeevashLongSkokan2021}
proved that an absolute constant $C\ge1$ suffices whenever
\begin{equation}\label{intro:kls}
 n\ge3,\qquad
 m\ge C\frac{\log_2 n}{\log_2\log_2 n}.
\end{equation}
This already contains every pair $m\ge n$ once $n$ is sufficiently large.
Indeed, the right-hand side of~\eqref{intro:kls} is $o(n)$.
For each of the finitely many smaller values of $n$, only finitely many
$m\ge n$ lie below that threshold. Thus their theorem leaves at most
finitely many pairs in the range of Theorem~\ref{thm:main}.
Our result completes that range. The finite calculation below comes from
an explicit structural reduction and does not require a numerical value
of the constant in~\eqref{intro:kls}.

\subsection{The proof}

Set $k=m-1$ and $a=n-1$. A failure of the upper bound gives a graph with
no cycle of order $k+1$ and independence number at most $a\le k$.
Holding $k$ fixed and minimizing $a$ produces a useful expansion property: every
nonempty independent set $I$ has at least $k|I|+1$ vertices in its closed
neighbourhood. In particular, the minimum degree is at least $k$.

The first structural step finds a clique of order
\[
 t\ge\max\{3,\lfloor k/2\rfloor\}.
\]
The proof uses distance layers in a rooted tree, followed by an explicit
path-rotation argument. Independent-set expansion and rooted distance
layers already appear in the cycle--clique argument of Erd\H{o}s,
Faudree, Rousseau and Schelp~\cite[Section~3]{ErdosFaudreeRousseauSchelp1978}.
We prove the particular quantitative statements needed here in
Sections~\ref{sec:reduction} and~\ref{sec:clique}.
The two smallest values $k=3,4$ are then settled directly.

For $k\ge5$, fix a maximum clique $Q$ of order $t$. Consider paths between
its vertices whose nonempty interiors lie outside $Q$ and are pairwise
disjoint, and whose endpoint pairs form disjoint chains. The total number
of interior vertices is called the \emph{amount} of the system. We maximize
this amount below $k+1-t$, breaking ties by minimizing the number of paths.
Certain improvements of this system would close to a cycle of order
exactly $k+1$; all such improvements are therefore forbidden.

These prohibitions make neighbourhoods of suitable path vertices
disjoint and anticomplete. Expansion then supplies independent sets
inside those neighbourhoods. A separate size lemma raises the available
independence bound from two to three when needed. For $t\ge9$, this
argument and a short numerical classification give more than $k$
independent vertices, a contradiction. Sections~\ref{sec:paths}--\ref{sec:terminal}
develop this part of the proof.

The remaining case has $3\le t\le8$, and the clique bound gives
$5\le k\le17$. A pattern records each chain by its sequence of
interior-vertex counts, up to chain reversal and permutation of components.
There are $3{,}099$ parameter-pattern instances across the $42$ admissible
pairs $(k,t)$. Section~\ref{sec:finite} proves the inference rules that
exclude them: path replacement, exact-cycle closure, neighbourhood
growth, independent-set packing, and contradiction under an added edge
or a path with one new interior vertex. The computation enumerates these
patterns, rather than graphs or edge colourings. Appendix~\ref{app:implementation}
describes two separate exact implementations and the accompanying
deduction traces. All structural arguments and all mathematical
inference rules are proved in the main text.
